% Fragmento público generado desde propietarios íntegros.
% Residencia: 52.3.2.
% Fuente: ../incoming/radion_monografia_20260827/manuscrito/sections/04_cincuenta_critica.tex:106-171
\subsubsection{Séparation typée des 2 objets de valeur \(1/2\)}

APP contient en outre un mode singulier avec
\[
s=\frac12,\qquad s^2=\frac14,\qquad
1-s^2=\frac34=\frac{90}{120},\qquad
s\sqrt{1-s^2}=\frac{\sqrt3}{4}.
\]
Ce \(s\) naît des angles principaux des feuillets somme et produit. La
couture critique \(\Re\rho=1/2\) naît d'une involution fonctionnelle. Que les deux
coordonnées valent \(1/2\) est une coïncidence scalaire dotée d'un contenu possible,
mais n'autorise pas à les identifier sans un carré typé.

La manière correcte de conserver la découverte est un diagramme d'égalisateurs :
\[
\begin{tikzpicture}[baseline=(current bounding box.center),node distance=14mm and 23mm,
 every node/.style={font=\small}]
\node[draw=RadBlue,rounded corners,fill=RadMist] (crit) {couture critique \(\Re s=1/2\)};
\node[draw=RadTeal,rounded corners,fill=RadCream,right=of crit] (face) {carte \(j_{100}\)};
\node[draw=RadCopper,rounded corners,fill=RadCream,right=of face] (theta) {phase \(\theta_2=50^\circ\)};
\node[draw=RadRose,rounded corners,fill=white,below=of face] (mode) {mode APP \(s=1/2\), \(1-s^2=90/120\)};
\draw[-{Latex},thick,RadBlue] (crit)--node[above]{\(j_{100}\)}(face);
\draw[-{Latex},thick,RadTeal] (face)--node[above]{égalisateur}(theta);
\draw[-{Latex},dashed,RadRose] (mode)--node[right,align=left]{même chiffre;\\domaine distinct}(face);
\end{tikzpicture}
\]
Le diagramme ne simule pas la flèche absente. Il montre ce qui est démontré, ce qui a
été formalisé et quelle égalité nécessite encore un transport supplémentaire si
l'on vise une identification opératoire globale.

\subsubsection{Compatibilité entre cercle spectral et mémoire hélicoïdale}

Le lecteur de Cayley--Pontryagin transporte une fibre spectrale au moyen de
\[
\widetilde U_{\Gamma_9^n}J\psi_\gamma
=\tau_\gamma^nJ\psi_\gamma.
\]
La publication circulaire est \(\tau_\gamma^n\in S^1\). Le relèvement
\[
(\tau_\gamma^n,n)\in S^1\times\Z
\]
conserve le nombre de retours. Sur le feuillet contractant apparaît
\[
\widetilde M_9^nJ\psi_\gamma
=\left(\frac59\right)^n\tau_\gamma^nJ\psi_\gamma.
\]
Nous avons donc trois images typées :
\[
\text{cercle de phase},\qquad
\text{hélice de mémoire},\qquad
\text{spirale contractante}.
\]

Un zéro n'est pas un tour, et un neuf n'est pas un zéro particulier de la fonction
zêta. Le zéro fixe un caractère \(\tau_\gamma\) qui persiste à travers les
tours ; le retour nonadique incrémente la mémoire. Le neuf fonctionne comme zéro
résiduel dans \(\Z/9\Z\), opération arithmétique distincte de l'énumération
des zéros spectraux.

\begin{narrativebox}
La droite critique devient un cercle lorsque nous oublions la hauteur linéaire et
conservons le caractère unitaire. Elle devient une hélice lorsque nous réinsérons
la mémoire du nombre de fois où le caractère a été transporté. La profondeur
n'est pas une troisième coordonnée décorative : c'est l'information que la projection
circulaire ne peut pas retenir.
\end{narrativebox}
